\section{The scalar Gaussian-message predecessor}
\label{app:scalar-prior}
This appendix is a reduction to prior work, not an additional novelty claim.
We state exactly which interpretation of \citet{Mahalanabis2012} is used.
Their Section 3.2, Theorem 43 approximates total conditional variance on a
Gaussian precision graph of treewidth $w$, with runtime
\begin{equation}
 \left(\frac{Nw\,\mathrm{cond}(\Lambda)}{\alpha}\right)^{O(w^2)}k^2
       +N2^{O(w^3)}
 \label{eq:gmrf-runtime}
\end{equation}
for relative variance error $1+\alpha$. Here $N$ counts Gaussian vertices,
not candidate times. The theorem as printed sums all unobserved variances
and permits all nodes as observations. We need only a single unobservable
target, and only designated observation candidates.

\subsection{The focused recurrence and a local-cost issue}
Their Definition 18, equation (37), assigns a conditional-variance cost to
the private variables of a subtree. Equations (39)--(40) combine those local
costs with child messages, while (41)--(43) pass separator precision by
observation and Schur complementation. Rounded counterparts are (44)--(48).
Restricting every local observed subset to designated candidates simply
removes choices; the proof's rounding step retains the same observed subset,
so every allowed optimal witness remains allowed. At the level of the covariance costs intended by Definition 18, replacing
each local trace by a nonnegative weighted sum of its private-coordinate variances
preserves relative-error comparisons: if
$e^{-a}M\preceq\widehat M\preceq e^aM$ with $M\succ0$, inversion gives
the same factors for each covariance diagonal and their nonnegative weighted
sum. The addition, observation, marginalization and rounding comparisons in
Observation 33 and Lemmas 41--42 then give the same induction. Zero costs
are compared by inequalities without division. This is an adaptation of the
proof, not the wording of the stated theorem.

One printed local-cost formula needs care. Taking an inverse after extracting
a private principal block generally differs from extracting that block after
inverting the full unobserved bag precision. For example,
\[
 M=\begin{pmatrix}2&1\\1&3\end{pmatrix},\qquad
 (M^{-1})_{11}=3/5\ne (M_{11})^{-1}=1/2.
\]
For the focused specialization we can avoid a nonzero cost in this disputed
form. Retain the target $\vartheta$ in every separator, direct messages
toward a singleton bag $\{\vartheta\}$ followed by an empty root, and give
all earlier private variables weight zero. The only positive cost is then
the target variance in that singleton bag, whose separator is empty and
where the two inverse/block orders coincide. Zero-weight local costs are
set to zero without evaluating an unnecessary inverse. Information messages
still use the stated Schur complements. Binary copies of bags retain
$\vartheta$ until the singleton and preserve this property. Thus the
focused construction does not depend on that erroneous local-cost identity.

\subsection{Rational normalization and regularization}
Take the scalar-observation class of Theorem~\ref{thm:trace-fptas}, with one
parameter ($p=1$) and cardinality as the only constraint. Cardinalities $k<0$
or $k>n$ are infeasible; for $k=0$, return the empty schedule. If all
sensitivities vanish, every feasible schedule is optimal. In the remaining
case, assume $1\le k\le n$ and at least one nonzero sensitivity $f_t$.
Choose positive powers of two $s,d_t$ with
\[
 s^2\le r_{\min}<4s^2,\qquad d_t^2\le r_t<4d_t^2.
\]
Define
\[
 U_t=X_t/s,\quad Z_t=Y_t/d_t,\quad c_t=s/d_t,
 \quad h=\max_t|f_t/d_t|>0,\quad g_t=f_t/(d_th).
\]
Then $|c_t|\le2$, $|g_t|\le1$, some $|g_t|=1$, and the normalized
observation-noise variances lie in $[1,4)$. The normalized parameter is
$\vartheta=h\theta$; its data information equals original data information
divided by $h^2$, which preserves rankings and relative ratios. The latent
covariance $K=\Cov(U)$ satisfies
\[
 |K_{ij}|\le4B_0\rho_0^{|i-j|},\qquad
 \norm K_2\le4B_0(1+\rho_0)/(1-\rho_0).
\]
For zero contraction the same formulas apply directly.

To allow zero original process variances, add independent variance $\zeta$
to each normalized initial/process innovation. Write $\mathsf A$ for the
strict lower shift with entries $a_t$. Its norm is at most $\rho_0$; thus
$T=(I-\mathsf A)^{-1}$ satisfies
\[
 \frac{\zeta}{(1+\rho_0)^2}I\preceq E:=\zeta TT^\top
       \preceq\frac{\zeta}{(1-\rho_0)^2}I.
\]
The lower bound follows from $\norm{T^{-1}}_2\le1+\rho_0$ and the upper
one from the Neumann series. Hence $K_+=K+E\succ0$ has polynomially
controlled spectrum. Let $R'$ be the normalized observation-error covariance
and $C=\diag(c_t)$. Since $R'\succeq I$ and $\norm C_2\le2$,
\[
 R'\preceq R'_+:=R'+CEC^\top
 \preceq R'+\frac{4\zeta}{(1-\rho_0)^2}I.
\]
Choose $\tau=\varepsilon/4$ and
$\zeta=\tau(1-\rho_0)^2/4$. Every selected principal covariance obeys
$R'_{SS}\preceq R'_{+,SS}\preceq(1+\tau)R'_{SS}$, so for
$I'(S)=g_S^\top(R'_{SS})^{-1}g_S$ and its regularized counterpart,
\begin{equation}
 I'(S)/(1+\tau)\le I'_+(S)\le I'(S).
 \label{eq:gmrf-regularization}
\end{equation}
All these changes are rational and have polynomial encoding length.

\subsection{Constant width and conversion of the objective}
Introduce an artificial independent target prior $\vartheta\sim N(0,1)$
and the model
\[
 U\sim N(0,K_+),\qquad Z_t=g_t\vartheta+c_tU_t+\nu_t,
 \qquad 1\le\Var(\nu_t)<4.
\]
Only the $Z_t$ are candidates; only $\vartheta$ has positive variance weight.
Then $v(S)=\Var(\vartheta\mid Z_S)=1/(1+I'_+(S))$. The Gaussian precision
graph has bags $\{\vartheta,U_t,U_{t+1},Z_t\}$ for $t<n$, followed by
$\{\vartheta,U_n,Z_n\}$ and the singleton target/empty root above.
These cover the latent-chain edges and each observation-density triangle,
and satisfy running intersection. Their width is at most three.

For conditioning, start from independent $(\vartheta,U,\nu)$ with covariance
$\diag(1,K_+,D)$, $I\preceq D\prec4I$. The shear to
$(\vartheta,U,Z)$ and its inverse have norms at most
$1+\sqrt{n+4}$, because the added bottom block $[g,C]$ has norm at most
$\sqrt{n+4}$. The joint covariance and inverse precision therefore have
condition number bounded by
\begin{equation}
 (1+\sqrt{n+4})^4
 \frac{\max\{4,\,4B_0(1+\rho_0)/(1-\rho_0)
                       +\zeta/(1-\rho_0)^2\}}
      {\min\{1,\,\zeta/(1+\rho_0)^2\}}
 =O_{\rho_0,B_0}(n^2/\varepsilon).
 \label{eq:gmrf-condition}
\end{equation}
In particular, normalization removes arbitrarily varying original noise and
sensitivity scales from this bound.

At an index with $|g_t|=1$, the regularized observation-error variance is
at most $D_0=4+16B_0+4\zeta/(1-\rho_0^2)$. A best nonempty schedule thus
has $I'_{+,*}\ge1/D_0$: start with that observation and use monotonicity to
pad to cardinality $k$. Apply the focused predecessor with
$\alpha=\varepsilon/[4(D_0+1)]$. Its variance ratio implies
\[
 I'_+(\widehat S)\ge\frac{I'_{+,*}-\alpha}{1+\alpha}
 \ge\frac{1-\alpha D_0}{1+\alpha}I'_{+,*}
 \ge(1-\varepsilon/4)I'_{+,*}.
\]
An at-most-$k$ output can be padded to exactly $k$, since observing more
variables cannot reduce information. Combining with
\eqref{eq:gmrf-regularization} yields
\[
 I'(\widehat S)\ge
 \frac{1-\varepsilon/4}{1+\varepsilon/4}\max_{|S|=k}I'(S)
 \ge(1-\varepsilon)\max_{|S|=k}I'(S).
\]
Multiplication by $h^2$ and addition of a common nonnegative scalar prior
preserve this guarantee. Equations~\eqref{eq:gmrf-runtime} and
\eqref{eq:gmrf-condition} make the predecessor's operation bound polynomial
for this normalized class.

This establishes the mathematical reduction in the predecessor's
spectral-rounding model. It does not silently supply a full rational
implementation of that rounding net. Our explicit bit-complexity result is
proved independently in Theorem~\ref{thm:trace-fptas}; the remaining numerical
model distinction is no basis for claiming a first scalar approximation
scheme. The reduction in this appendix uses cardinality alone and does not
justify padding under spacing or arbitrary side constraints.

\section{Proof of the represented-matroid cover}
\label{app:matroid-proof}
We prove Theorem~\ref{thm:matroid-spectral-set} using the owner-count
variant of the exact profile construction. The topic-22 Lean package
formalizes this variant, including its original-input execution and
polynomial bit-work bound. The accompanying repository records the exact
claim coverage, independent reviews, and targeted verification commands.

If $q=\rank A=0$, the only base is empty and is returned, regardless of
the prior. Otherwise select $q$ independent rows of $A$ and discard the
redundant rows, preserving its column matroid. Apply
Lemma~\ref{lem:spectral-normalization}, retaining its range and magnitude
tests. For a rank-$r$ trial, write $E'$ for the retained elements and $F$
for the distinct forced owners. Only original bases contained in $E'$
and containing $F$ may be returned. Such a base exists only if
$F\subseteq E'$, $A_F$ is independent, and $\rank A_{E'}=q$.
These conditions can be tested directly or through the exact support test
below. A smaller base of a rank-deficient restriction is never accepted.

\subsection{Forcing owners without contraction}
Retain the original $q$ rows, and zero columns outside $E'$; call this
matrix $A_{\rm ret}$. Give each element the marker
$b_e=\mathbf1\{e\in F\}$. Since a base contains each element at most once,
\begin{equation}
 \sum_{e\in B}b_e=|F|
 \quad\Longleftrightarrow\quad F\subseteq B.
 \label{eq:matroid-owner-marker}
\end{equation}
One extra integer profile coordinate therefore enforces all owners.
This replaces rational contraction without changing the accepted bases.
Write $q'=q-|F|$ for their common number of optional elements. If $q'=0$,
every accepted base is exactly $F$ and has zero approximation error; the
same support procedure may still be used.

\subsection{Exact profile detection and witness recovery}
Put $h=\eta/(rq)$ and $z_{e,ij}=\lfloor A_{e,ij}/h\rfloor$ for $i\le j$.
There are $d=r(r+1)/2$ information coordinates. The prior is unrounded.
Forced elements may be rounded with the other retained elements because
their common contribution cancels between accepted bases. Set
$L_* =\lceil4p/h\rceil$ and
$w_{e,ij}=z_{e,ij}+L_*\in\{0,\ldots,2L_*\}$.
Every original base has $q$ elements, so subtracting $qL_*\mathbf1$
recovers its signed profile. Subtracting the shared forced profile leaves
the optional profile of $q'$ elements.

Over the rational field define
\begin{equation}
\begin{split}
 g(y,t)&=\det\!\left(A_{\rm ret}
               \diag_e(y^{w_e}t^{b_e})A_{\rm ret}^{\top}\right)\\
 &=\sum_{\substack{B\text{ original base}\\B\subseteq E'}}
       \det(A_B)^2 y^{\sum_{e\in B}w_e}t^{|B\cap F|}.
\end{split}
 \label{eq:matroid-profile-polynomial}
\end{equation}
This is the Cauchy--Binet identity used by
\citet[Lemmas 4.3--4.4]{Berstein2008}, with one additional integer
coordinate. Expanding over $q$-column subsets gives zero for dependent
subsets and the displayed squared minor for every base. Thus every
coefficient is nonnegative and is positive exactly when its profile is
attained. Retain only marker degree $|F|$, which enforces
\eqref{eq:matroid-owner-marker}. Characteristic-zero arithmetic is essential:
a positive coefficient may vanish modulo a prime. A zero residue alone
cannot certify infeasibility.

The information degrees are at most $D_*=2qL_*$ and the marker degree is
at most $|F|\le q$. Use the uniform grid
$\{1,\ldots,\overline D+1\}^{d+1}$, where
$\overline D=\max(D_*,q)$. Since $L_*\ge1$ in these trials,
$\overline D=D_*$. A sharper rectangular grid is also possible.
Exact tensor Lagrange interpolation recovers the requested coefficient:
for each univariate node $j$, form
\[
 \ell_j(X)=\prod_{\substack{1\le s\le\overline D+1\\s\ne j}}
                   \frac{X-s}{j-s},
\]
and sum each grid value multiplied by the product of the corresponding
coefficients of the $\ell_j$. This is an explicit inverse of the tensor
Vandermonde system. The coefficient-vector products construct the
univariate polynomials with polynomial work; they do not enumerate all
subsets of factors. It is equivalent to the established exact
interpolation method, whose original presentation uses a moment curve.

For every positive owner-complete profile, recover one original base by
one-pass deletion. Start with the retained columns and delete a column
exactly when the requested coefficient remains positive after its removal.
Keep the original $q$ rows in all tests: rank loss makes the polynomial
zero. Never replace it by a smaller row basis. Feasibility of the target
persists, and a failed deletion cannot become feasible after further
columns are removed. An extra surviving column outside one target-profile
base would contradict its failed deletion. Thus at most $m$ tests leave
exactly an original base with that profile. No lifting is needed.

\subsection{Coverage, zero information, and bit bounds}
Fix an arbitrary target base $B$ with information rank $r>0$. Its
maximum-volume trial from Lemma~\ref{lem:spectral-normalization} retains
all of $B$. The forced owners form a subset of $B$, and the retained set
has original rank $q$. The marker-complete profile of $B$ therefore has a
positive coefficient. Let $\widehat B$ be its recovered representative.
If $q'=0$, both bases equal $F$. Otherwise their optional signed profiles
are equal after cancelling $F$, and each optional entry residual lies in
$[0,q'h)$. The transformed information entries consequently differ by
less than $q'h\le qh$. Their symmetric difference has norm at most
$rqh=\eta$. The forced identity floor and rectangular reconstruction give
\eqref{eq:spectral-cover} and exact range equality. For information rank
zero, require $Q_0=0$ and find an original-rank base using only zero-matrix
elements. A nonzero prior excludes this branch.

There are at most $\binom Mr$ rank-$r$ trials. Fixing the marker at $|F|$
and subtracting the common forced profile leaves at most
$R_*=(2q'L_*+1)^d$ positive profiles. The extra interpolation coordinate
therefore increases work but not the number of returned bases. In particular,
\[
 |\mathcal C|\le1+\sum_{r=1}^p\binom Mr
 \bigl[2q\lceil4prq/\eta\rceil+1\bigr]^{r(r+1)/2}.
\]
The initial one permits the zero-information representative; $q'=0$
contributes at most one base in its trial.

The variable matrix rank requires polynomial determinant evaluation,
not a permutation expansion. The implementation clears denominators and
uses a cached division-free Bird recurrence on the integer matrix
\citep{Bird2011}.
It stores each preceding matrix once; an order-$q$ determinant takes
$O(q^4)$ ring operations. Bounds on the integer iterates and exact final
division by the common denominator raised to $q$ give polynomial bit
work, even when $q$ grows. Rational row selection and information
normalization have polynomial bit bounds as well.

At a grid point the monomial has bit length
$O((dL_*+1)\log(\overline D+1))$. The evaluated matrix, determinant,
univariate Lagrange coefficients, tensor terms, and finite sums all have
polynomial encoding lengths. The grid has polynomial size for fixed
$d+1$. Recomputing its values for each requested coefficient and each
deletion test still gives polynomial work in input length and $1/\eta$
for fixed $p$. Basis output length is polynomial too. This completes the
mathematical Turing-model argument without a floating zero test,
independence-oracle conversion, or feasibility relaxation.
The formal cost model charges schoolbook arithmetic, finite scans, storage,
and copies, including rejected trials and labels. It excludes construction
of proof and cost-observer transcripts and makes no Lean wall-clock claim.

\subsection{Exact witnesses for the representation restrictions}
The squared-minor argument is specific to a single represented matroid.
For two representations the mixed determinant contains products of different
minors, whose signs may differ. With row matrices $A_1=(1,1)$ and
$A_2=(1,-1)$, and identical zero profiles, both singleton sets are common
bases but $\det(A_1A_2^\top)=0$. Thus direct substitution loses feasible
profiles. A separate randomized profile algorithm for represented matroid
intersection is established by \citet[Section 3]{Berstein2010}; element
indeterminates and random substitution address that cancellation. We do not
claim its guarantee as part of our deterministic theorem. Likewise,
\[
 \det\begin{pmatrix}1&1&0\\1&0&1\\0&1&1\end{pmatrix}=-2
\]
is nonzero over $\Q$ but zero over $\mathbb F_2$, so reinterpreting a
finite-field representation can change the bases.

Projection does not remove attainability requirements either. The rank-one
uniform matroid on two elements with scalar information $1$ and $3$ has
projected information polytope $[1,3]$. Its integer point $2$ is not attained
by any base. Integer optimization over that interval would therefore solve
a different problem. Finally, the independence-oracle result of
\citet[Theorem 1.1]{Berstein2008} assumes a fixed number of distinct weight
values. The number of our rounded labels grows with input size and
$1/\eta$; substituting them in that oracle theorem does not retain a fixed
polynomial exponent. These examples explain limitations of the proof routes,
not hardness of every possible broader spectral-cover algorithm.
